\documentclass[11pt,a4paper]{article}
\usepackage[UTF8,scheme=plain]{ctex}
\usepackage[margin=2cm]{geometry}
\usepackage{amsmath,booktabs,enumitem,hyperref}
\usepackage[labelsep=quad,font=small]{caption}
\usepackage{float}
\usepackage{titlesec}
\titleformat*{\section}{\large\bfseries}
\titlespacing*{\section}{0pt}{1.4ex plus .3ex}{.7ex}
\titlespacing*{\subsection}{0pt}{1ex plus .2ex}{.5ex}
\setlist{nosep,leftmargin=1.6em}
\linespread{1.18}
\setlength{\parskip}{0.2em}
\hypersetup{colorlinks=true,urlcolor=blue}
\AtBeginDocument{\renewcommand{\tablename}{表}}

\title{\bfseries 跨细胞类型的药物响应预测及其误差来源\\[0.3em]
\large 统计学习大作业开题报告}
\author{盛其然 \quad 洪华言 \quad 应硕丞 \quad 樊宇轩}
\date{2026 年 10 月}

\begin{document}
\vspace*{-2.5em}
{\let\newpage\relax\maketitle}
\vspace*{-1.5em}

\section{问题背景}

如果对每一种 ``药物 $\times$ 细胞类型'' 组合都做单细胞实验，成本很高。
NeurIPS 2023 单细胞扰动竞赛因此提出：\textbf{能否用某些细胞类型里测到的药物响应，预测同一药物在另一种细胞类型里的响应？}
实验用人外周血单个核细胞（PBMC）加入 144 种小分子药物，24 小时后测单细胞基因表达。
NK 和三类 T 细胞几乎测了所有药物，B 细胞和髓系细胞（Myeloid）只公开了 15 种药物，
需要预测它们其余 255 个组合的响应。

竞赛前排方案大多已开源，主流做法是把 ``药物'' 和 ``细胞类型'' 编码为特征（如目标均值编码），
在降维后的输出空间上训练神经网络或梯度提升等多输出模型，部分方案还引入了药物化学结构等外部信息。
这些方案的模型框架大同小异，但较少分析\textbf{模型为什么有效、误差从何而来}，这正是本项目希望补充的部分。

\section{具体研究问题}

核心是\textbf{如何更准确地预测}（问题一）。
问题二至四是围绕数据本身的若干小实验，从标签、噪声、生物学机制三个角度寻找影响预测的因素：
一方面为改进模型提供方向，另一方面帮助解释模型为什么有效或失效。

\subsection*{问题一（核心）：如何用其他细胞类型的响应，更准确地预测目标细胞类型？}
具体包括四个方面。
\begin{itemize}
\item \textbf{响应的结构}：药物响应在细胞类型之间有多少共性？
  18{,}211 维的响应中，大多数基因对大多数药物几乎没有显著变化，真正承载信号的可能只是少数基因。
  因此除了 ``低维'' 结构，我们更关心响应是否\textbf{稀疏}，以及能否先筛出有响应的基因再建模。
  预实验显示响应\textbf{极度重尾}：平方和最大的 5\% 的行占全部平方和的 81\%。
\item \textbf{简单模型与复杂模型}：线性模型、核方法与神经网络差多少，复杂模型的增益来自哪里。
\item \textbf{模型的可解释性}：神经网络本身难以解释，但我们可以在其基础上加入有明确含义的改进模块，
  再通过消融与敏感性分析说明每个模块为什么带来提升。
\item \textbf{跨细胞类型的泛化}：模型在不同目标细胞类型上的表现是否一致，表现差的细胞类型和药物有什么共同特点。
\end{itemize}

\subsection*{问题二：标签是怎样生成的？这对预测意味着什么？}
标签衡量的是\textbf{显著性}而非变化幅度：同样的药效，参与拟合的样本越多，显著性越大。
主办方计算 B/Myeloid 训练标签时，可能只用了 17 种训练药物的孔（``小模型''），也可能连测试药物的孔一起用（``大模型''）。
二者可用训练孔中读数全为 0 的基因区分：小模型下这类基因的 17 行标签必然完全相同；大模型下（在 T/NK 上检验）只有 35\%--51\% 相同。
实际 B/Myeloid 标签中这类基因\textbf{全部}相同（B 139/139，Myeloid 30/30），说明\textbf{训练标签来自小模型}（自由度约 130，大模型约 380），
而测试标签很可能来自大模型。因此训练标签的数值系统性偏小，模型学到的幅度也会偏小。
我们将量化这一差异，并研究如何在训练和验证中加以考虑。

\subsection*{问题三：实验噪声给预测设了多高的上限？}
标签含有任何模型都预测不出的随机噪声，因此预测误差存在下限。
官方测试标签不公开，我们无法直接得知模型在测试集上 ``最好能做到多好''，但可以借助原始单细胞数据估计噪声的大小。
预实验中，我们分别算出 3 位供体各自的表达变化：NK 和 CD4$^+$ 中供体间相关系数的中位数只有 0.04 和 0.12，最强的 10\% 的药物也仅 0.26 和 0.37，
即\textbf{大多数药物的响应在供体之间几乎重复不出来}。
我们要估计误差下限，回答现有模型离它多远、剩余误差来自方向、幅度还是噪声，从而判断改进的方向。

\subsection*{问题四：细胞类型特异性响应能否由药物靶点的表达解释？}
预测依赖借用其他细胞类型中同一药物的响应；若药物只在某一种细胞类型中起强作用，所有模型都会失效。
一个可能的解释是：药物通过结合靶点或受体蛋白起作用，只有表达该靶点的细胞才会响应。
我们打算从公开数据库整理药物靶点，用对照样本算出靶点在各细胞类型中的表达，检验这一解释是否普遍成立，
再把 ``靶点在目标细胞中是否表达'' 作为新特征加入模型。
这既回答了 ``药物本身的信息能否帮助预测''，也为模型提供生物学上的解释。

\section{数据来源}

数据来自 \href{https://www.kaggle.com/competitions/open-problems-single-cell-perturbations/data}{Kaggle 竞赛页面}。
\begin{itemize}
\item \textbf{主数据} \texttt{de\_train.parquet}：614 行 ``细胞类型 $\times$ 药物'' 组合，每行是 18{,}211 个基因的
  $\text{sign}(\log\text{FC})\times(-\log_{10}p)$。B 和 Myeloid 只有 15 种药物（下称\textbf{锚点药物}）及 2 种对照。
\item \textbf{辅助数据}：官方公开的训练部分原始单细胞计数 \texttt{adata\_train.parquet} 及样本信息（3 位供体），用于问题二至四。
\item \textbf{外部数据}：药物靶点从 LINCS、ChEMBL 等公开数据库整理；竞赛前排的开源方案作为参考。
\end{itemize}
不使用任何测试集标签或测试集原始数据。

\section{初步建模思路}

以下为初步设想，具体方法将根据实验结果调整。

\textbf{问题一：预测模型。}
\begin{itemize}
\item \textbf{基线与线性模型}：全零预测、已知药物均值等朴素基线，以及岭回归、核岭回归等线性或核方法，作为比较的基准。
\item \textbf{输出降维与稀疏化}：输出维度（18{,}211）远大于样本量（614），需要先压缩输出。
  我们将比较两种思路：用截断 SVD 等方法得到低维表示，或先筛出在多数药物上有显著响应的少数基因，只对这些基因建模。
\item \textbf{多任务学习}：把每个已知细胞类型的预测看作一个任务，各任务共享药物层面的特征，
  再把学到的共享结构迁移到 B/Myeloid。
\item \textbf{神经网络}：在已有的神经网络模型（参考开源的前排方案）基础上，加入有明确含义的改进模块，
  如针对重尾的幅度模块、问题二的标签尺度校准、问题四的靶点特征，以及基因本身的信息（如基因在各细胞类型中的基础表达）。
\item \textbf{消融与敏感性分析}：逐个去掉模块或遮盖部分输入，观察误差变化，判断各部分的重要性；
  对输入施加小幅扰动（如 5\%--10\%），检验预测是否稳定。
\end{itemize}

\textbf{问题二至四：数据分析。}
\begin{itemize}
\item \textbf{问题二}：用原始计数复现 Limma 流程，模拟小模型与大模型，量化两者的差异，并检验据此校准预测幅度能否改进预测。
\item \textbf{问题三}：按供体拆分原始数据，估计标签噪声，得到误差下限；把各模型误差拆成方向、幅度、噪声三部分。
\item \textbf{问题四}：整理靶点注释并计算靶点在各细胞类型中的表达，检验其与响应特异性的关系，再作为特征加入模型。
\end{itemize}

\section{评价方案}

评价遵循两条原则：一是验证集要\textbf{模拟真实任务}，即目标细胞类型只知道少数锚点药物，因此不用随机划分；
二是\textbf{测试数据不参与训练和调参}。
初步设想两种方案：
\textbf{方案 A} 轮流把一个数据完整的 T/NK 细胞类型当作 ``新细胞类型''，只保留其锚点药物，其余药物作为测试，超参数在其余细胞类型上用同样方式选取；
\textbf{方案 B} 在 B/Myeloid 的锚点药物上做留一验证，直接针对目标细胞类型，但样本较少，作为补充。
评价指标以竞赛指标 MRRMSE 为主，并辅以置信区间等不确定性度量；方案 A 还可以分别报告各目标细胞类型上的表现，用于回答问题一中的泛化问题。
具体方案将随研究进展调整，例如问题二若证实训练与测试标签的生成方式不同，验证集也需相应修正。

\section{可能困难}

\begin{enumerate}
\item \textbf{输出维度远大于样本量}，且大部分基因接近噪声：通过降维或基因筛选压缩输出，并比较不同压缩方式的影响。
\item \textbf{神经网络可解释性差}：不直接解释网络本身，而是通过消融与敏感性分析解释各改进模块的作用。
\item \textbf{供体只有 3 位}，噪声估计不稳定：尝试多种拆分方式并报告敏感性。
\item \textbf{锚点药物只有 15 种}，B/Myeloid 上样本少：主要在 T/NK 上检验，再到 B/Myeloid 上验证。
\item \textbf{原始单细胞数据量大}，读取和处理较慢：先汇总为伪 bulk 数据，并利用课题组的 GPU 服务器计算。
\end{enumerate}

\section{初步分工与进度}

初步分工如下，后续可根据进展和个人兴趣灵活调整：
应硕丞负责问题一（预测模型、消融与敏感性分析）；盛其然负责统筹与问题二；洪华言负责问题三；樊宇轩负责问题四。
重要结果由不同成员独立复核（交叉验证），代码统一整理在 GitHub 仓库中，保证可复现。

进度安排：
\begin{itemize}
\item \textbf{10 月}：整理并共享已有模型和代码，统一训练与验证流程；开始问题二至四的数据分析。
\item \textbf{11 月上旬}（期中考试期间）：以数据分析为主，完成中期汇报。
\item \textbf{11 月中旬起}：在已有模型上进行改进、消融与敏感性实验，把问题二至四的结论接入模型。
\item \textbf{12 月}：统一评估，整理代码，撰写最终报告。
\end{itemize}

\end{document}
